
\documentclass{article}
\usepackage{colortbl}
\usepackage{makecell}
\usepackage{multirow}
\usepackage{supertabular}

\begin{document}

\newcounter{utterance}

\centering \large Interaction Transcript for game `cladder', experiment `full\_v1.5\_default', episode 916 with refineonly{-}Qwen3.5{-}9B.
\vspace{24pt}

{ \footnotesize  \setcounter{utterance}{1}
\setlength{\tabcolsep}{0pt}
\begin{supertabular}{c@{$\;$}|p{.15\linewidth}@{}p{.15\linewidth}p{.15\linewidth}p{.15\linewidth}p{.15\linewidth}p{.15\linewidth}}
    \# & \multicolumn{2}{c}{Player} && \multicolumn{2}{c}{Game Master} \\
    \hline

    \theutterance \stepcounter{utterance}  
    & & & \multicolumn{4}{p{0.6\linewidth}}{
        \cellcolor[rgb]{0.9,0.9,0.9}{
            \makecell[{{p{\linewidth}}}]{
                \texttt{\tiny{[P1$\langle$GM]}}
                \texttt{You are an expert in causal inference. The following question is not a typical commonsense query, but rather a meticulously designed question created by a professor specializing in causal inference, intended to assess the students' mastery of the course content.} \\
\texttt{Imagine a self{-}contained, hypothetical world with only the following conditions, and without any unmentioned factors or causal relationships: Smoking has a direct effect on effort and college admission. Effort has a direct effect on college admission. For nonsmokers and are lazy, the probability of college admission is 47\%. For nonsmokers and are hard{-}working, the probability of college admission is 88\%. For smokers and are lazy, the probability of college admission is 23\%. For smokers and are hard{-}working, the probability of college admission is 54\%. For nonsmokers, the probability of being hard{-}working is 40\%. For smokers, the probability of being hard{-}working is 61\%. If we disregard the mediation effect through effort, would smoking negatively affect college admission?} \\
\texttt{Start your answer with 'Yes' or 'No', followed by additional reasoning or evidence to support your explanation.} \\
            }
        }
    }
     \\ \\

    \theutterance \stepcounter{utterance}  
    & \multicolumn{4}{p{0.6\linewidth}}{
        \cellcolor[rgb]{0.561,0.737,0.561}{
            \makecell[{{p{\linewidth}}}]{
                \texttt{\tiny{[P1$\rangle$GM]}}
                \texttt{No} \\
\\ 
\texttt{**Reasoning:**} \\
\\ 
\texttt{To determine if smoking negatively affects college admission when disregarding the mediation effect through effort, we must calculate the **Total Effect** of smoking on admission. The total effect is the difference in the probability of admission between the counterfactual states where an individual smokes versus does not smoke, averaging over the distribution of the mediator (effort).} \\
\\ 
\texttt{Mathematically, the total effect \$TE\$ is defined as:} \\
\texttt{\$\$TE = P(Y\_\{smoker\} = 1) {-} P(Y\_\{nonsmoker\} = 1)\$\$} \\
\\ 
\texttt{Where \$Y\_\{smoker\}\$ is the outcome if the individual smokes, and \$Y\_\{nonsmoker\}\$ is the outcome if the individual does not smoke. We calculate these probabilities by summing the conditional probabilities of admission given effort, weighted by the probability of being hard{-}working (or lazy) for each smoking status.} \\
\\ 
\texttt{**Step 1: Calculate the probability of admission for a Nonsmoker (\$P(Y\_\{nonsmoker\} = 1)\$)**} \\
\texttt{*   Probability of being hard{-}working given nonsmoker: \$P(E=1|S=0) = 0.40\$} \\
\texttt{*   Probability of being lazy given nonsmoker: \$P(E=0|S=0) = 1 {-} 0.40 = 0.60\$} \\
\texttt{*   Admission probability given Nonsmoker \& Hard{-}working: \$P(Y=1|S=0, E=1) = 0.88\$} \\
\texttt{*   Admission probability given Nonsmoker \& Lazy: \$P(Y=1|S=0, E=0) = 0.47\$} \\
\\ 
\texttt{\$\$P(Y\_\{nonsmoker\} = 1) = (0.40 \textbackslash{}times 0.88) + (0.60 \textbackslash{}times 0.47)\$\$} \\
\texttt{\$\$P(Y\_\{nonsmoker\} = 1) = 0.352 + 0.282 = 0.634\$\$} \\
\\ 
\texttt{**Step 2: Calculate the probability of admission for a Smoker (\$P(Y\_\{smoker\} = 1)\$)**} \\
\texttt{*   Probability of being hard{-}working given smoker: \$P(E=1|S=1) = 0.61\$} \\
\texttt{*   Probability of being lazy given smoker: \$P(E=0|S=1) = 1 {-} 0.61 = 0.39\$} \\
\texttt{*   Admission probability given Smoker \& Hard{-}working: \$P(Y=1|S=1, E=1) = 0.54\$} \\
\texttt{*   Admission probability given Smoker \& Lazy: \$P(Y=1|S=1, E=0) = 0.23\$} \\
\\ 
\texttt{\$\$P(Y\_\{smoker\} = 1} \\
            }
        }
    }
    & & \\ \\

    \theutterance \stepcounter{utterance}  
    & & & \multicolumn{2}{p{0.3\linewidth}}{
        \cellcolor[rgb]{0.9,0.9,0.9}{
            \makecell[{{p{\linewidth}}}]{
                \texttt{\tiny{[GM$|$GM]}}
                \texttt{No} \\
            }
        }
    }
    & & \\ \\

    \theutterance \stepcounter{utterance}  
    & & & \multicolumn{2}{p{0.3\linewidth}}{
        \cellcolor[rgb]{0.9,0.9,0.9}{
            \makecell[{{p{\linewidth}}}]{
                \texttt{\tiny{[GM$|$GM]}}
                \texttt{yes} \\
            }
        }
    }
    & & \\ \\

    \theutterance \stepcounter{utterance}  
    & & & \multicolumn{2}{p{0.3\linewidth}}{
        \cellcolor[rgb]{0.9,0.9,0.9}{
            \makecell[{{p{\linewidth}}}]{
                \texttt{\tiny{[GM$|$GM]}}
                \texttt{game\_result = LOSE} \\
            }
        }
    }
    & & \\ \\

\end{supertabular}
}

\end{document}
